\documentclass[a4paper,12pt]{report}
\addtolength{\oddsidemargin}{-1.cm}
\addtolength{\textwidth}{2cm}
\addtolength{\topmargin}{-2cm}
\addtolength{\textheight}{3.5cm}
\newcommand{\HRule}{\rule{\linewidth}{0.5mm}}
\makeindex

\usepackage{longtable}
\usepackage[pdftex]{graphicx}
\usepackage{makeidx}
\usepackage{hyperref}
\hypersetup{
    colorlinks=true,
    linkcolor=blue,
    filecolor=magenta,      
    urlcolor=cyan,
}


% define the title
\author{Group3_a}
\title{ Assignment 1 Report}
\begin{document}
\setlength{\parskip}{6pt}

% generates the title
\begin{titlepage}

\begin{center}
% Upper part of the page       
\includegraphics[width=1\textwidth]{./University_of_Pretoria_Logo.PNG}\\[0.4cm]    
\textsc{\LARGE Department of Computer Science}\\[1.5cm]
\textsc{\Large COS 301 - Software Engineering}\\[0.5cm]
% Title
\HRule \\[0.4cm]
{ \huge \bfseries COS 301 - Mini Project}\\[0.4cm]
\HRule \\[0.4cm]
% Author and supervisor
\begin{minipage}{0.4\textwidth}
\begin{flushleft} \large
\emph{Author:}\\
Sean {Hill}
\end{flushleft}
\end{minipage}
\begin{minipage}{0.4\textwidth}
\begin{flushright} \large
\emph{Student number:} \\
u12221458
\end{flushright}
\end{minipage}
\begin{minipage}{0.4\textwidth}
\begin{flushleft} \large
Liz {Joseph}
\end{flushleft}
\end{minipage}
\begin{minipage}{0.4\textwidth}
\begin{flushright} \large
\emph{} \\
u10075268
\end{flushright}
\end{minipage}
\begin{minipage}{0.4\textwidth}
\begin{flushleft} \large
Jessica {Lessev}
\end{flushleft}
\end{minipage}
\begin{minipage}{0.4\textwidth}
\begin{flushright} \large
\emph{} \\
u13049136
\end{flushright}
\end{minipage}
\begin{minipage}{0.4\textwidth}
\begin{flushleft} \large
Kale-ab {Tessera}
\end{flushleft}
\end{minipage}
\begin{minipage}{0.4\textwidth}
\begin{flushright} \large
\emph{} \\
u13048423
\end{flushright}
\end{minipage}
\begin{minipage}{0.4\textwidth}
\begin{flushleft} \large
Sphelele {Malo}
\end{flushleft}
\end{minipage}
\begin{minipage}{0.4\textwidth}
\begin{flushright} \large
\emph{} \\
u12247040
\end{flushright}
\end{minipage}
\begin{minipage}{0.4\textwidth}
\begin{flushleft} \large
Diana {Obo}
\end{flushleft}
\end{minipage}
\begin{minipage}{0.4\textwidth}
\begin{flushright} \large
\emph{} \\
u13134885
\end{flushright}
\end{minipage}
\begin{minipage}{0.4\textwidth}
\begin{flushleft} \large
Armand {Piterse}
\end{flushleft}
\end{minipage}
\begin{minipage}{0.4\textwidth}
\begin{flushright} \large
\emph{} \\
u12167844
\end{flushright}
\end{minipage}
\vfill
% Bottom of the page
{\large \today}
\end{center}
\end{titlepage}
\footnotesize
%%
%%  University of Pretoria - Declaration of Originality
%%  Version: S4726/09 (amended 2010-07-23)
%%

\newpage

\thispagestyle{empty}
{
\renewcommand{\baselinestretch}{1}
\sffamily
\begin{center}
\textbf{\Large DECLARATION OF ORIGINALITY}
\end{center}
\begin{center}
\textbf{\large UNIVERSITY OF PRETORIA}
\end{center}

The University of Pretoria places great emphasis upon integrity and
ethical conduct in the preparation of all written work submitted for
academic evaluation.

While academic staff teach you about referencing techniques and how to
avoid plagiarism, you too have a responsibility in this regard. If you
are at any stage uncertain as to what is required, you should speak to
your lecturer before any written work is submitted.

You are guilty of plagiarism if you copy something from another
author's work (e.g. a book, an article or a website) without
acknowledging the source and pass it off as your own. In effect you
are stealing something that belongs to someone else. This is not only
the case when you copy work word-for-word (verbatim), but also when
you submit someone else's work in a slightly altered form (paraphrase)
or use a line of argument without acknowledging it. You are not
allowed to use work previously produced by another student. You are
also not allowed to let anybody copy your work with the intention of
passing if off as his/her work.

Students who commit plagiarism will not be given any credit for
plagiarised work. The matter may also be referred to the Disciplinary
Committee (Students) for a ruling. Plagiarism is regarded as a serious
contravention of the University's rules and can lead to expulsion from
the University.

The declaration which follows must accompany all written work
submitted while you are a student of the University of Pretoria. No
written work will be accepted unless the declaration has been
completed and attached.

\begin{center}
\begin{tabular}{ll}
                        &                             \\
 Full names of students: & \makebox[3.5in]{\hrulefill} \\ \\
 Student numbers:        & \makebox[3.5in]{\hrulefill} \\ \\
 Topic of work:         & \makebox[3.5in]{\hrulefill} \\ \\
\end{tabular}
\end{center}

\textbf{Declaration}
\begin{enumerate}
\item I understand what plagiarism is and am aware of the University's
  policy in this regard.
\item I declare that this assignment report is my own original work.
  Where other people's work has been used (either from a printed source,
  Internet or any other source), this has been properly acknowledged and
  referenced in accordance with departmental requirements.
\item I have not used work previously produced by another student or any
  other person to hand in as my own.
\item I have not allowed, and will not allow, anyone to copy my work with
  the intention of passing it off as his or her own work.
\end{enumerate}

\vspace*{1cm}

SIGNATURES: \makebox[3in]{\hrulefill} \quad DATE: \makebox[1.5in]{\hrulefill}
}

\newpage



\normalsize

\renewcommand{\thesection}{\arabic{section}}
\newpage
\begin{center}
\textsc{\LARGE Architectural Requirements}\\[1.5cm]
\textsc{\Large Buzz Space link}\\[0.5cm]
For further references see \href{https://github.com/COS301-Group-3-A-Phase2/Part2-Mini-Project.git}{gitHub}.
\today
\end{center}
\tableofcontents{}

\newpage

% You are welcome to edit these
% They change according to how you structure your work.
\section{Change Log}

\section{Introduction}
Designing and implementing a system as the one in question is a vast amount of work and requires a lot of planning and thinking. All the functional requirements needed for the system were covered in the previous phase. In this phase all the non-functional requirements will be covered. These are all categorized as architectural requirements.  \\
\subsection{Purpose}
The main purpose of this document is to decompose the system at hand and fully describe the decomposition in terms of subsystem responsibilities, the mapping of the subsystem to hardware and the dependencies among the subsystems. Policy decisions that need to be made is also described in this document. This included decisions such as access control, data storage and control flow.\\\\
The document will cover all the non-functional aspects of the requirements including:\\
•	the architectural scope comprising the responsibilities,\\
•	the quality requirements for the software system,\\
•	the access and integration requirements, and\\
•	the architectural constrains. \\

\subsection{Project Scope}
The scope of this phase is understanding the functional requirements covered in the previous phase and using them to build and derive the architectural requirements needed for the implementation of the system.\\

\section{Access Channel Requirements}
\subsection{Platforms}

\section{Quality Requirements}
\subsection{Performance}
\subsection{Reliability}
\subsection{Scalability}
\subsection{Security}
\subsection{Flexibility}
\subsection{Maintainability}
\subsection{Auditability}
\subsection{Monitorability}
\subsection{Cost}
\subsection{Usability}

\section{Integration Requirements}
\subsection{Integration Channel}
\subsection{Protocols}
\subsection{API Specifications}
\subsection{Quality Requirements}

\section{Architecture Constraints}
\subsection{Technologies}
\subsection{Architectural Patterns/Frameworks}

\section{Policies}

\section{Traceability Matrix}

\section{References}
PIETERSE, V. 2015. COS301: Mini project Buzz. In: Lecture notes issued online. University of Pretoria. Pretoria, South Africa

\section{Glossary}

\end{document}

\bibliography{myrefs}{} 
\bibliographystyle{ieeetr}
\end{document}
